% GENERATED FILE. Edit canonical lessons, then run npm run generate:books.
% canonical-source-hash: fnv1a64:a096195a7d755081
% canonical-lessons: ZH-C442-duian, ZH-C442-he, ZH-W442-he-river, ZH-C442-shangyou, ZH-C442-zhangshui, ZH-C442-jiang

\chapter{The River}
\label{ch:zh-the-river}
\hlchaptermodality{442}

\begin{chapteropening}
\textbf{By the end of this chapter:} \emph{I can talk about the far bank, a river, its upper reaches, a river in flood and a great river.}

Complete the last lesson of chapter 442: I can talk about the far bank, a river, its upper reaches, a river in flood and a great river.
\end{chapteropening}

\section[duì'àn]{duì'àn — the far bank; the opposite bank}

\label{lesson:ZH-C442-duian}



\begin{quote}
Before the new one: say the Mandarin for to light a fire, then the Mandarin for made by nature.
\end{quote}

\subsection*{You'll want to know: duì'àn}
\emph{duì'àn} — \textquotedblleft{}the far bank; the opposite bank\textquotedblright{}.

\begin{quote}
\emph{wǒmen zuò chuán qù duì'àn} — We're taking a boat across to the other side.
\end{quote}

Its characters come later; for now it is a word for the ear.

\begin{cousinweb}
\emph{duì'àn} is built from \zh{对} 'facing' and \emph{àn} 'a bank', both inherited from Middle Chinese, \emph{twojH} and \emph{nganH} (Guangyun, in Baxter's transcription), which glosses the second as the high edge of the water.
\end{cousinweb}

\begin{grammarlens}[title={What you've built}]
One more word for getting around.
\end{grammarlens}

\subsection*{Your turn}
\begin{itemize}
\raggedright
  \item \emph{Say it:} \emph{duì'àn}
  \item \emph{Say it:} \emph{duì'àn}, once more
  \item \emph{Say it:} \emph{wǒmen zuò chuán qù duì'àn} — We're taking a boat across to the other side.
  \item \emph{Recall:} say the Mandarin for to light a fire, then the Mandarin for made by nature, then say \emph{duì'àn} again
\end{itemize}

\begin{tcolorbox}[breakable,colback=teal!4,colframe=teal!35!black,title={Before you move on}]
What does duì'àn mean? (\textquotedblleft{}The far bank; the opposite bank\textquotedblright{}.) Say it once more.
\end{tcolorbox}
\section[\texorpdfstring{\zh{河} (hé)}{河 (hé)}]{\zh{河} (hé) — river}

\label{lesson:ZH-C442-he}



\begin{quote}
Before the new one: say the Mandarin for made by nature, then the Mandarin for the far bank.
\end{quote}

\subsection*{You'll want to know: \zh{河}}
\textbf{\zh{河}} — \emph{hé} — \textquotedblleft{}river\textquotedblright{}.

\begin{quote}
\emph{cūnzi qiánbian yǒu yì tiáo hé} — There's a river in front of the village.
\end{quote}

The new character, in pieces: \zh{河} — eight strokes in two pieces: on the left the water piece of \zh{洗}, three strokes; on the right \zh{可} kě, five strokes. The next lesson puts it in your hand.

\begin{cousinweb}
\textbf{\zh{河}} (\emph{hé}) is inherited from Middle Chinese \emph{ha} (Guangyun, in Baxter's transcription), where it is the name of one river, the one said to rise in the far west and run to the sea; the sense 'any river' is usually explained as a widening of that name.
\end{cousinweb}

\begin{grammarlens}[title={What you've built}]
One more word for getting around.
\end{grammarlens}

\subsection*{Your turn}
\begin{itemize}
\raggedright
  \item \emph{Say it:} \emph{hé}
  \item \emph{Say it:} \emph{hé}, once more
  \item \emph{Say it:} \emph{cūnzi qiánbian yǒu yì tiáo hé} — There's a river in front of the village.
  \item \emph{Recall:} say the Mandarin for made by nature, then the Mandarin for the far bank, then say \emph{hé} again
\end{itemize}

\begin{tcolorbox}[breakable,colback=teal!4,colframe=teal!35!black,title={Before you move on}]
What does \zh{河} mean? (\textquotedblleft{}River\textquotedblright{}.) Say it once more.
\end{tcolorbox}
\section[\texorpdfstring{\zh{河} (hé)}{河 (hé)}]{\zh{河} — the water piece, then \zh{可}}

\label{lesson:ZH-W442-he-river}



\begin{quote}
The word first, then the character inside it.

\begin{itemize}
\raggedright
  \item \emph{Recall:} say \emph{hé} — \textbf{R1}, one lesson back, before it decays
\end{itemize}

\begin{quote}
\zh{洗} \zh{可}
\end{quote}

\emph{Read it:} each one aloud, with nothing said first — \emph{xǐ}, \emph{kě}
\end{quote}

\subsection*{Script — the water piece, then \zh{可}}
\hlblockfigure{\detokenize{figures/ZH-W442-he-river-filmstrip.pdf}}{How \zh{河} is written, stroke by stroke}

\begin{quote}
\zh{河}
\end{quote}

\textbf{Two pieces, left before right.} Count the pieces, not the strokes.

\textbf{The left piece, the water piece, three strokes}, as in \zh{洗}: two dots, then a stroke that rises to the right.

\textbf{The right piece, \zh{可} kě, five strokes}, as you wrote it: a long bar; under it \zh{口}, three strokes, set to the left; then a long vertical down the right side that hooks to the left at the foot.

3 + 5 = eight.

\subsection*{Your turn}
\begin{itemize}
\raggedright
  \item \emph{Write it:} \zh{河}, all eight, with the model in front of you
  \item \emph{Cover:} write it again from memory, and count the pieces, not the strokes
\end{itemize}

\begin{tcolorbox}[breakable,colback=teal!4,colframe=teal!35!black,title={Before you move on}]
How many strokes in \zh{河}? (\textbf{Eight}.) What is the right piece? (\textbf{\zh{可} kě}.) What is the last stroke? (\textbf{The long vertical with a hook}.)

Source: \href{https://raw.githubusercontent.com/chanind/hanzi-writer-data/68d10a4b21150cae5e1ebbd223eed289cf32d90c/data/\%E6\%B2\%B3.json}{Hanzi Writer Data \zh{河}, pinned stroke paths and medians}.
\end{tcolorbox}
\section[\texorpdfstring{\zh{上游} (shàngyóu)}{上游 (shàngyóu)}]{\zh{上游} (shàngyóu) — upper reaches (of a river); upstream}

\label{lesson:ZH-C442-shangyou}



\begin{quote}
Before the new one: say the Mandarin for the far bank, then the Mandarin for a river.
\end{quote}

\subsection*{You'll want to know: \zh{上游}}
\textbf{\zh{上游}} — \emph{shàngyóu} — \textquotedblleft{}upper reaches (of a river); upstream\textquotedblright{}.

\begin{quote}
\textbf{\zh{我们坐船去河的上游}} \emph{wǒmen zuò chuán qù hé de shàngyóu} — We're taking a boat up the river.
\end{quote}

\begin{cousinweb}
\textbf{\zh{上游}} (\emph{shàngyóu}) is built from \zh{上} 'above' and \zh{游} 'to float', both inherited from Middle Chinese, \emph{dzyangH} and \emph{yuw} (Guangyun, in Baxter's transcription); how the second came to name a stretch of a river is not settled.
\end{cousinweb}

\begin{grammarlens}[title={What you've built}]
One more word for getting around.
\end{grammarlens}

\subsection*{Your turn}
\begin{itemize}
\raggedright
  \item \emph{Say it:} \emph{shàngyóu}
  \item \emph{Say it:} \emph{shàngyóu}, once more
  \item \emph{Say it:} \emph{wǒmen zuò chuán qù hé de shàngyóu} — We're taking a boat up the river.
  \item \emph{Recall:} say the Mandarin for the far bank, then the Mandarin for a river, then say \emph{shàngyóu} again
\end{itemize}

\begin{tcolorbox}[breakable,colback=teal!4,colframe=teal!35!black,title={Before you move on}]
What does \zh{上游} mean? (\textquotedblleft{}Upper reaches (of a river); upstream\textquotedblright{}.) Say it once more.
\end{tcolorbox}
\section[zhǎngshuǐ]{zhǎngshuǐ — to rise after rain (of a river)}

\label{lesson:ZH-C442-zhangshui}



\begin{quote}
Before the new one: say the Mandarin for a river, then the Mandarin for upstream.
\end{quote}

\subsection*{You'll want to know: zhǎngshuǐ}
\emph{zhǎngshuǐ} — \textquotedblleft{}to rise after rain (of a river)\textquotedblright{}.

\begin{quote}
\emph{zhè jǐ tiān xiàyǔ hé lǐ zhǎngshuǐ le} — It has rained these last few days, and the river has risen.
\end{quote}

Its characters come later; for now it is a word for the ear.

Two third tones meet in this word: zhǎng is said with a rising tone before shuǐ, and the spelling keeps both marks.

\begin{cousinweb}
\emph{zhǎngshuǐ} is built from \emph{zhǎng} 'to swell' and \zh{水} 'water', both inherited from Middle Chinese, \emph{trjang} and \emph{sywijX} (Guangyun, in Baxter's transcription), which glosses the first as a great flood of water; today's falling-rising tone is irregular beside it.
\end{cousinweb}

\begin{grammarlens}[title={What you've built}]
One more word for getting around.
\end{grammarlens}

\subsection*{Your turn}
\begin{itemize}
\raggedright
  \item \emph{Say it:} \emph{zhǎngshuǐ}
  \item \emph{Say it:} \emph{zhǎngshuǐ}, once more
  \item \emph{Say it:} \emph{zhè jǐ tiān xiàyǔ hé lǐ zhǎngshuǐ le} — It has rained these last few days, and the river has risen.
  \item \emph{Recall:} say the Mandarin for a river, then the Mandarin for upstream, then say \emph{zhǎngshuǐ} again
\end{itemize}

\begin{tcolorbox}[breakable,colback=teal!4,colframe=teal!35!black,title={Before you move on}]
What does zhǎngshuǐ mean? (\textquotedblleft{}To rise after rain (of a river)\textquotedblright{}.) Say it once more.
\end{tcolorbox}
\section[jiāng]{jiāng — a great river; big river}

\label{lesson:ZH-C442-jiang}



\begin{quote}
Before the new one: say the Mandarin for upstream, then the Mandarin for to rise after rain.
\end{quote}

\subsection*{You'll want to know: jiāng}
\emph{jiāng} — \textquotedblleft{}a great river; big river\textquotedblright{}.

\begin{quote}
\emph{zhè tiáo jiāng bǐ nà tiáo hé cháng} — This great river is longer than that one.
\end{quote}

Its characters come later; for now it is a word for the ear.

\begin{cousinweb}
\emph{jiāng} is inherited from Middle Chinese \emph{kaewng} (Guangyun, in Baxter's transcription), which glosses it with the pair 'rivers and seas' and lists the great rivers of the south under it.
\end{cousinweb}

\begin{grammarlens}[title={What you've built}]
One more word for getting around.
\end{grammarlens}

\subsection*{Your turn}
\begin{itemize}
\raggedright
  \item \emph{Say it:} \emph{jiāng}
  \item \emph{Say it:} \emph{jiāng}, once more
  \item \emph{Say it:} \emph{zhè tiáo jiāng bǐ nà tiáo hé cháng} — This great river is longer than that one.
  \item \emph{Recall:} say the Mandarin for upstream, then the Mandarin for to rise after rain, then say \emph{jiāng} again
\end{itemize}

\begin{tcolorbox}[breakable,colback=teal!4,colframe=teal!35!black,title={Before you move on}]
What does jiāng mean? (\textquotedblleft{}A great river; big river\textquotedblright{}.) Say it once more.
\end{tcolorbox}
